\talkslide{defender-loss}{Learning the defender: online DPO}{Timing: provisional}{For the defender, the candidates are two responses to the same prompt.}{
\dpolossview{1}
}
\speak{The diagram shows a schematic reference-relative score profile, with the preferred and rejected candidates as individual points. Language-model outputs are discrete and the candidate ordering is arbitrary. The connecting curve is only a visual guide, not a probability density or measured curve. Scores can be positive or negative; the schematic has no numeric vertical scale. The vertical arrows depict raising the preferred candidate’s score and lowering the rejected candidate’s score. The exact objective compares their reference-relative log probabilities, not their absolute probabilities alone. The score is beta times the log probability relative to the reference model, not a scalar score emitted by the judge. The expectation averages this loss over freshly collected preference pairs. Winner and loser are denoted w and l. The compact notation suppresses the fixed conditioning variables inside the score. The defender follows the same online DPO construction we used for the attacker, now with responses rather than attacks as the candidates. The judge chooses the response that is more compliant for a benign seed or more constructively deflecting for a harmful seed. The loss raises the winner's likelihood relative to the loser and a reference policy. The displayed equation follows the paper; the released configuration also supports an EMA reference. The difference is where the prompts and preference pairs come from in the joint game. Source detail: AdvGame, Eq. (8); $w$ and $l$ denote winner and loser. Reference choice revisited on slide \ref{deck-rollouts}.}
